\section{Problem setup}
\label{sec:setup}

\paragraph{Prompted style transfer with a frozen task model.}
Let $x$ be a source sentence with an undesired attribute (here, a negative Yelp
review). A policy $\pol(\cdot \mid \lmb)$ emits a short discrete
\emph{prompt} $z = (z_1,\dots,z_L)$ of fixed length $L=5$ tokens. The prompt is
inserted into a fixed template together with the source,
\begin{equation}
  \texttt{template}(z, x) \;=\; \texttt{\{z\}\ "\{x\}"\ "} ,
\end{equation}
and a \emph{frozen}, inference-only task language model $\lm$ samples a
rewrite $y \sim \lm(\cdot \mid \texttt{template}(z,x))$. Only $\theta$ is
trained; $\lm$ receives no gradient. This is the prompt-optimization regime: the
action space is the discrete prompt, and the environment --- the task model plus
the scorers --- is a black box.

\paragraph{Two objectives.}
Each rewrite is scored on two axes, both on a $0$--$100$ scale:
\begin{itemize}\itemsep2pt
  \item \emph{content} $c(x,y) \in [0,100]$, a BERTScore-style similarity
        between the rewrite and its source, measuring how much of the original
        meaning survives;
  \item \emph{sentiment} $s(y) \in [0,100]$, the probability assigned to the
        target attribute by a fine-tuned BERT classifier.
\end{itemize}
The two are in genuine tension: copying $x$ verbatim maximizes $c$ and leaves
$s$ at the source value, while emitting a generic positive sentence maximizes
$s$ and destroys $c$.

\paragraph{A $\lmb$-indexed \emph{constrained} reward.}
Rather than a linear scalarization, the objectives are combined as a
constraint. For $\lmb \in [0,1]$, define
\begin{equation}
  \Rlam(z,x) \;=\;
  \begin{cases}
    s(y), & c(x,y) \;\ge\; 100\lmb, \\[2pt]
    0.01\,\bigl(c(x,y) - 100\lmb\bigr), & \text{otherwise,}
  \end{cases}
  \label{eq:reward}
\end{equation}
so $\lmb$ sets a \emph{content floor} at $100\lmb$: reward equals the sentiment
score wherever the floor is met, and otherwise degrades to a small negative
penalty proportional to the shortfall. Sweeping $\lmb$ from $0$ (no content
requirement --- pure sentiment) to $1$ (near-verbatim content required) traces
the tradeoff curve. The penalty branch is scaled by $0.01$ so that violation is
weakly informative rather than a flat rejection, which keeps the objective
differentiable-in-expectation and gives the policy a gradient to climb back
into the feasible region.

Because $\lm$ is stochastic, the reward used for learning is a Monte Carlo
average over $N = 50$ task-model samples per (prompt, source) pair,
\begin{equation}
  \hat{r}_{\lmb}(z,x) \;=\; \frac{1}{N}\sum_{n=1}^{N} \Rlam(z, x; y_n),
  \qquad y_n \sim \lm(\cdot \mid \texttt{template}(z,x)),
  \label{eq:mcreward}
\end{equation}
which is what every algorithm below consumes as its scalar reward.

\paragraph{Objective: one policy for the whole curve.}
During training $\lmb$ is drawn afresh per source, $\lmb \sim \mathcal{U}[0,1]$,
and supplied to the policy as a conditioning input. A single set of parameters
$\theta$ must therefore serve every point on the curve, so that at test time the
frontier is traced by varying the conditioning input rather than by retraining:
\begin{equation}
  \max_{\theta}\;
  \mathbb{E}_{\lmb \sim \mathcal{U}[0,1]}\;
  \mathbb{E}_{x}\;
  \mathbb{E}_{z \sim \pol(\cdot \mid \lmb)}\bigl[\hat{r}_{\lmb}(z,x)\bigr].
  \label{eq:objective}
\end{equation}
This conditioning is what makes the problem multi-objective rather than a set of
independent single-objective runs, and it is the reason entropy collapse is
especially damaging here: a policy that collapses onto one prompt becomes
$\lmb$-independent and degenerates to a single point instead of a curve.
